\documentclass[lettersize,journal]{IEEEtran}

% ============================================================
% PACKAGES
% ============================================================
\usepackage{amsmath,amssymb,amsfonts,amsthm}
\usepackage{algorithm}
\usepackage{algpseudocode}
\usepackage{graphicx}
\usepackage{cite}
\usepackage{url}
\usepackage{xcolor}
\usepackage{booktabs}
\usepackage{multirow}
\usepackage{bm}
\usepackage[caption=false,font=footnotesize]{subfig}
\usepackage{balance}
\usepackage{tikz}
\usetikzlibrary{arrows.meta,positioning,calc,decorations.pathmorphing,
  shapes.geometric,fit,backgrounds,patterns,fadings,shadows.blur}
\usepackage{pgfplots}
\pgfplotsset{compat=1.18}

\hyphenation{op-tical net-works semi-conduc-tor Kolmo-go-rov}

% ============================================================
% THEOREM ENVIRONMENTS
% ============================================================
\newtheorem{corollary}{Corollary}
\newtheorem{lemma}{Lemma}
\newtheorem{remark}{Remark}
\newtheorem{definition}{Definition}
\newtheorem{proposition}{Proposition}
\newtheorem{assumption}{Assumption}

% ============================================================
% CUSTOM COMMANDS
% ============================================================
\DeclareMathOperator{\erfc}{erfc}
\DeclareMathOperator{\diag}{diag}
\DeclareMathOperator*{\argmin}{arg\,min}
\DeclareMathOperator*{\argmax}{arg\,max}
\newcommand{\vect}[1]{\mathbf{#1}}
\newcommand{\mat}[1]{\mathbf{#1}}
\newcommand{\set}[1]{\mathcal{#1}}
\newcommand{\norm}[1]{\left\|#1\right\|}
\newcommand{\abs}[1]{\left|#1\right|}
\newcommand{\R}{\mathbb{R}}
\newcommand{\E}{\mathbb{E}}

% ============================================================
\usepackage[active,tightpage]{preview}
\setlength\PreviewBorder{0.6pt}
\begin{document}
\begin{preview}
\begin{tikzpicture}[
  cell/.style={circle, draw=black!70, fill=#1, minimum size=11pt,
               inner sep=0pt, font=\tiny\bfseries},
  cell/.default={blue!15},
  mol/.style={circle, fill=orange!80, minimum size=2.2pt, inner sep=0pt},
  edge/.style={-,gray!50, line width=0.6pt},
  diffedge/.style={-{Stealth[length=3pt,width=2.5pt]}, blue!60!black,
                   line width=0.8pt,
                   decorate, decoration={snake, amplitude=0.8pt,
                   segment length=4pt, post length=2.5pt}},
  lbl/.style={font=\scriptsize, text=black!80},
  eqlbl/.style={font=\footnotesize, text=black},
  block/.style={rectangle, rounded corners=3pt, draw=black!60,
                fill=#1, minimum height=18pt, minimum width=40pt,
                font=\scriptsize, align=center, line width=0.5pt},
  block/.default={white},
  bigblock/.style={rectangle, rounded corners=5pt, draw=black!40,
                   fill=#1, minimum height=28pt, font=\small,
                   align=center, line width=0.7pt, text width=52pt},
  bigblock/.default={white},
]

% ============================================================
% PANEL (a): Biological setting — cells in tissue
% ============================================================
\begin{scope}[shift={(0,0)}]
  \node[lbl, font=\small\bfseries] at (2.9,3.65) {(a) Biological setting};

  % Background tissue
  \fill[blue!3, rounded corners=6pt] (-0.3,-0.5) rectangle (6.8,3.35);
  \draw[blue!20, rounded corners=6pt, line width=0.6pt]
       (-0.3,-0.5) rectangle (6.8,3.35);
  \node[lbl, blue!40, font=\scriptsize] at (5.85,3.05) {aqueous medium};

  % Cell positions (randomized look)
  \coordinate (c1) at (0.7, 2.4);
  \coordinate (c2) at (2.1, 2.7);
  \coordinate (c3) at (1.3, 1.2);
  \coordinate (c4) at (3.4, 2.1);
  \coordinate (c5) at (2.7, 0.7);
  \coordinate (c6) at (4.5, 2.6);
  \coordinate (c7) at (4.0, 1.0);
  \coordinate (c8) at (5.2, 1.7);
  \coordinate (c9) at (0.5, 0.3);
  \coordinate (c10) at (3.8, 0.0);

  % Edges (communication links by proximity)
  \foreach \i/\j in {1/2, 1/3, 2/3, 2/4, 3/5, 4/5, 4/6, 5/7, 6/8, 7/8,
                     7/10, 3/9, 9/5}{
    \draw[edge] (c\i) -- (c\j);
  }

  % Diffusion arrows (selected links)
  \draw[diffedge] (c2) -- node[above,lbl,pos=0.45]{$w_{ij}$} (c4);
  \draw[diffedge] (c4) -- (c6);

  % Molecules floating
  \foreach \x/\y in {1.5/2.15, 1.75/2.4, 2.6/2.3, 2.45/1.85,
                     3.0/1.55, 3.8/1.65, 4.2/2.2, 0.9/0.8,
                     3.2/0.45, 4.7/2.1}{
    \node[mol] at (\x,\y) {};
  }

  % Cells
  \node[cell=green!20]  at (c1)  {$1$};
  \node[cell=green!35]  at (c2)  {$i$};
  \node[cell=green!15]  at (c3)  {$3$};
  \node[cell=red!20]    at (c4)  {$j$};
  \node[cell=green!25]  at (c5)  {$5$};
  \node[cell=green!20]  at (c6)  {$6$};
  \node[cell=green!10]  at (c7)  {$7$};
  \node[cell=green!15]  at (c8)  {$8$};
  \node[cell=green!10]  at (c9)  {$9$};
  \node[cell=green!10]  at (c10) {\!10\!};

  % Distance annotation
  \draw[|<->|, black!60, line width=0.4pt]
       ([yshift=-4pt]c2.south) -- node[below, lbl]{$d_{ij}$}
       ([yshift=-4pt]c4.south);

  % Cell detail inset (arrow from node j, label above-left)
  \draw[black!30, line width=0.4pt, ->] (c4.north) -- ++(0.3, 0.7)
       node[above, lbl]{radius $r_r$,\; $\mathbf{c}_j(t) \in \mathbb{R}^M$};
\end{scope}

% ============================================================
% PANEL (b): Graph abstraction and channel model
% ============================================================
\begin{scope}[shift={(7.8,0)}]
  \node[lbl, font=\small\bfseries] at (3.2,3.65) {(b) Communication graph $\mathcal{G}$};

  \fill[gray!4, rounded corners=6pt] (-0.3,-0.5) rectangle (6.8,3.35);
  \draw[gray!25, rounded corners=6pt, line width=0.6pt]
       (-0.3,-0.5) rectangle (6.8,3.35);

  % Graph nodes
  \coordinate (n1) at (0.6,2.5);
  \coordinate (n2) at (2.0,2.8);
  \coordinate (n3) at (1.1,1.3);
  \coordinate (n4) at (3.2,2.2);
  \coordinate (n5) at (2.5,0.8);

  % Edges with weights
  \draw[edge, line width=1.44pt, blue!29] (n1) -- (n2);
  \draw[edge, line width=2.00pt, blue!35] (n1) -- (n3);
  \draw[edge, line width=0.50pt, blue!20] (n2) -- (n3);
  \draw[edge, line width=1.44pt, blue!29] (n2) -- (n4);
  \draw[edge, line width=1.17pt, blue!27] (n3) -- (n5);
  \draw[edge, line width=0.86pt, blue!24] (n4) -- (n5);
  \node[cell=blue!25, minimum size=14pt] at (n1) {$1$};
  \node[cell=blue!35, minimum size=14pt] at (n2) {$i$};
  \node[cell=blue!20, minimum size=14pt] at (n3) {$3$};
  \node[cell=blue!40, minimum size=14pt] at (n4) {$j$};
  \node[cell=blue!25, minimum size=14pt] at (n5) {$5$};

  % Weight label
  \node[lbl, blue!70, rotate=12] at (2.55,2.65) {$w_{ij}$};

  % CIR plot inset (shifted right)
  \begin{scope}[shift={(4.0,0.0)}]
    \draw[black!40, ->, line width=0.4pt] (0,0) -- (2.3,0)
         node[right, lbl]{$t$};
    \draw[black!40, ->, line width=0.4pt] (0,0) -- (0,2.0)
         node[above, lbl]{$h_{ij}(t)$};
    % CIR curve (heavy tail)
    \draw[red!70!black, line width=1.0pt, smooth]
         plot coordinates {
           (0.08,0.0) (0.15,0.5) (0.25,1.6) (0.35,1.45) (0.5,1.0)
           (0.7,0.65) (0.9,0.45) (1.1,0.33) (1.4,0.22) (1.7,0.16)
           (2.0,0.12) (2.2,0.10)};
    % Tail annotation
    \node[lbl, red!60!black, font=\tiny] at (1.85,0.45) {$\sim t^{-3/2}$};
    % Integration area (w_ij)
    \fill[blue!10, opacity=0.6]
         (0.08,0) -- plot coordinates {
           (0.08,0.0) (0.15,0.5) (0.25,1.6) (0.35,1.45) (0.5,1.0)
           (0.7,0.65) (0.9,0.45)} -- (0.9,0) -- cycle;
    \draw[blue!50, line width=0.4pt, dashed] (0.9,0) -- (0.9,0.45);
    \node[lbl, blue!70, font=\tiny] at (0.5,-0.2) {$w_{ij}{=}\!\int_0^\tau\!h\,dt$};
    \node[lbl, black!50, font=\tiny] at (0.9,-0.35) {$\tau$};
  \end{scope}

\end{scope}

\end{tikzpicture}
\end{preview}
\end{document}
